\section{The weighted annulus estimate}\label{sec:annulus}

Theorem~\ref{thm:circuit-input} follows from a mixed-cylinder crossing
estimate and the renormalization argument of \cite[Appendix~A]{GL2025}.
We prove the required seed estimate using both signs of the first-spin
crossing and both orientations of the cell lattice.  A winding
construction converts a seed of either orientation to the one required
by the cylinder convention.

\subsection{Cuts and comparisons on a cylinder}

Write $Y$ and $Z$ for the two classes of elementary triangle centers,
with $Y$ the class used in Section~\ref{sec:discrete}, and put $q=x^2$.
An occupation of class $T\in\{Y,Z\}$ uses the cell table
\eqref{discrete:cell-table} on triangles of that class.
For a periodic rows-cylinder of height $H$, let $P_T$ denote its mixed
law: $B=+$ on its top face layer, $B=-$ on its bottom face layer,
$W$ free, with factors $q$ on changing interior $T$ cells and no factors
on boundary $T$ cells. Thus the class specifies the boundary convention
as well as the augmentation. All contours and test regions below have
positive clearance; lattice rounding changes their positions by $O(1)$.

\begin{lemma}[Ambient extension]\label{annulus:ambient}
Take a finite triangulation of the sphere containing a periodic triangular
lattice cylinder. Choose any edge-disjoint family of elementary triangles
in that cylinder, make each a cell of parameter $q=x^2$, and make each
remaining edge a cell of parameter $x$, or of parameter $1$ on a specified
set of edges. Use coherent spin pairs and the cell table
\eqref{discrete:cell-table}. Proposition~\ref{prop:history-order} holds
for this law, including deterministic $W$-equality edges and its stated
first-spin histories. The selected triangles may have either orientation.

A simple authorized winding contour in the cylindrical portion, with no
$W$-equality edge crossing it, has the cut property of
Lemma~\ref{lem:cell-cut}. In particular, prescribing a
lower open-minus and an upper open-plus contour of class $T$ realizes
$P_T$ between those contours, when they bound the rows-cylinder.
\end{lemma}
\begin{proof}
Cap a longer periodic cylinder at each end by a triangulated disk;
triangular fans suffice. Every mod-two cycle on this sphere is generated
by triangular face boundaries. Consequently, for fixed $B=b$, a possible
$W$ disagreement set satisfies the same triangle equations as in the proof
of Proposition~\ref{prop:history-order}. On a mixed-$b$ triangle its
single homogeneous edge cannot belong to that set. Restricting to
homogeneous-plus edges or to homogeneous-minus edges therefore separately
preserves every triangle equation. Each restriction integrates to a spin
configuration, and the wall factorization \eqref{discrete:wall-factor}
holds, with its same factor $1/2$.

For an unobserved triangle cell, $r_y(b)$ is $1$ or $x^2$ according as
$b$ is constant or mixed; for an edge cell it is $1$ or that edge's
parameter. On an unobserved homogeneous cell the corresponding allowed
wall-edge parameter is $x$, or the specified single-edge parameter;
on an observed open cell it is $1$. Equality edges have parameter zero
in the wall partition sums. The cell inequality for $r_y$ and the
intersection/union inequalities for the two allowed-edge vectors are
unchanged. The Ising partition-sum inequality
\eqref{discrete:Z-lattice} thus proves the same two-law inequality.
Restoring statuses uses only disjointness of the cells' edge sets.
This proves association and conditional order, including parameters $1$.

For the cut assertion, include among the selected triangles all cells
responsible for authorizing the contour. Each crossing cell has its first
spins prescribed and contributes a constant, exactly as in
Lemma~\ref{lem:cell-cut}; the remaining factors separate. This proof is
local and applies to a winding contour, which separates the capped sphere.
The equality constraints do not cross the contour by hypothesis.
Two such cuts isolate the law between them. The excluded boundary-cell
factors are precisely the constants removed by the cuts. No extra winding
constraint is imposed: an annular configuration satisfying the cut-layer
pins extends coherently to either cap by taking the prescribed first spin
constant there, and the exterior partition factors are independent of the
interior second spins.
\end{proof}

In applications of Lemma~\ref{annulus:ambient}, triangles of opposite
orientations are selected only in separated regions. Unused augmentations
are marginalized, and all remaining edges are single-edge cells. We do not
assert joint association for overlapping $Y$ and $Z$ augmentations.

\begin{lemma}[One complementary-color cut]\label{annulus:one-white}
The comparison of Lemma~\ref{annulus:ambient} for $B$ and signed
$B$ statuses remains valid with one prescribed open $W$ winding contour
of fixed sign, provided that its cells are disjoint from the selected
$B$ cells. Histories still observe no other $W$ labels.
\end{lemma}
\begin{proof}
An open $W$ triangle with prescribed sign contributes the constant $q$
for every assignment of its $B$ spins. Represent its prescription by
single-edge parameters $1$ on its three edges and by $W$ equality along
those edges, then specify the common $W$ label. The incident-site union
of the contour cells is connected: consecutive cells along the contour
share a site. Before that label is specified, global $W$ flip pairs the
two possible labels with equal weight, for every $B$ and every permitted
$B$-status history. Specifying its sign therefore multiplies each such
marginal weight by $1/2$. It does not change the order comparison.
Apply Lemma~\ref{annulus:ambient} with these deterministic equality edges.
\end{proof}

\begin{lemma}[Boundary slides]\label{annulus:slide}
Let $\mu_{a,b}$ be the law between a lower open-minus $B$ contour $a$
and an upper open-plus $B$ contour $b$. If the contours are ordered
from bottom to top as
\[
 a_1<a_0<\operatorname{supp}(A)<b_1<b_0,
\]
where $A$ is an increasing $B$ and signed-status event, then
\begin{equation}\label{annulus:slide-cuts}
 \mu_{a_0,b_0}(A)\le \mu_{a_1,b_1}(A).
\end{equation}
Old and new contours may use different classes, with edge-disjoint cells.
In particular, writing $\tau_d A_T$ for the upward translate by $d$ of
an increasing class-$T$ test, for $S,T\in\{Y,Z\}$ we have
\begin{equation}\label{annulus:slide-tests}
 P_S(A_T)\le P_T(\tau_d A_T),
\end{equation}
whenever the downward-translated class-$T$ boundary contours have the
required clearance from the original class-$S$ contours and the test.

Above a fixed open $W$ contour $\gamma$, let $\nu_{\gamma,b}$ denote
the law with upper open-plus $B$ contour $b$. For
$\gamma<\operatorname{supp}(A)<b_0<b_1$, the reverse slide gives
\begin{equation}\label{annulus:slide-white}
 \nu_{\gamma,b_1}(A)\le\nu_{\gamma,b_0}(A).
\end{equation}
\end{lemma}
\begin{proof}
Work on a common sufficiently long capped cylinder. Prescribe both
$a_0$ and $a_1$. By the cut property, the law above $a_0$ is unaffected
by the additional prescription at $a_1$. Releasing the old open-minus
prescription at $a_0$ increases the probability of $A$, by
Lemma~\ref{annulus:ambient}. Next prescribe $b_1$, which increases that
probability, and discard the farther contour $b_0$ by the cut property.
This proves \eqref{annulus:slide-cuts}. Translate the new cylinder back
up by $d$ to obtain \eqref{annulus:slide-tests}; the lattice translation
preserves its class $T$.

For \eqref{annulus:slide-white}, prescribe both $b_0$ and $b_1$ and then
release the nearer positive prescription at $b_0$. This decreases the
probability of $A$, by Lemma~\ref{annulus:one-white}. The distant
prescription did not affect the original law below $b_0$, by the cut
property. All simultaneous prescriptions used here are feasible: one
may take the complementary spin constant; when the fixed $W$ contour
is present, take that constant to be its prescribed sign.
\end{proof}

\begin{lemma}[Exterior search for a winding cut]\label{annulus:search}
In a cylindrical band, search from below for an authorized winding contour
of a specified class, color, and sign. The search can be arranged so that,
conditional on each successful complete record, the law above the selected
contour is the corresponding cut law. Existence of a winding with fixed
clearance inside the search band implies success.
\end{lemma}
\begin{proof}
Use the face-invasion construction after Lemma~\ref{discrete:online}.
Start from the lower face layer and invade across unauthorized honeycomb
edges, querying a responsible cell together with all its spins of the
searched color. Restrict queries to the search band; declare failure if
the invasion reaches its upper edge. An authorized winding with clearance
prevents that failure. The boundary against the uninvaded component toward
the upper end contains a separating authorized winding contour. Retain
such a simple contour, discarding contractible boundary components.
Trivalence prevents a four-edge crossing at a boundary vertex.

Every queried cell touches the invaded face set. If it is not favorable,
all its incident edges are unauthorized and the invasion reaches all its
incident faces. Thus a queried cell meeting the retained side must be a
favorable contour cell. The only first spins revealed there are the
prescribed cut-layer spins; complementary-spin labels are never queried.
The record's remaining restrictions lie below the contour. The product
factorization of Lemma~\ref{annulus:ambient}, with the colors interchanged
when necessary, now identifies the conditional law above it. We condition
on the actual successful search record, not on an additional event in the
retained region. A bounded enlargement of the band accommodates conversion
between triangular-cell windings and honeycomb contours.
\end{proof}

\begin{lemma}[Signed band alternative]\label{annulus:duality}
In a cylindrical band augmented in either class $T$, at least one of the
following events occurs: a vertical occupied $B+$ crossing, a vertical
occupied $B-$ crossing, an occupied $W+$ winding, or an occupied $W-$
winding. Each alternative yields one of a fixed finite family of local
horizontal or vertical crossing tests when the band width and circumference
are fixed multiples of its height. The number of tests is independent of
the lattice scale.
\end{lemma}
\begin{proof}
Couple the two color occupations using opposite tests on each uniform, as
in Definition~3.5 of \cite{GL2025}. Since $q\ge1/2$, every cell belongs
to at least one occupation. If there is no vertical $B$-occupied path,
triangular-site path/circuit duality on the cylinder gives a winding of
cells outside $B$ occupation, hence inside $W$ occupation. Adjacent occupied
cells of a given color share a site, so their spin signs agree. Both a
vertical path and a winding therefore have one constant sign, proving the
four-way alternative.

For the last assertion, lift a path to the flat strip cover and let $m$
be the band height. If a vertical path's horizontal extent is at most
$5m$, one of finitely many width-$6m$ windows contains it. Otherwise it
has a horizontal subcrossing of width greater than $4m$. A winding also
has such a horizontal subcrossing. Choose horizontal centers on a net of
spacing at most $m/4$ and leave fixed slack for lattice rounding. Taking
subpaths yields the local crossing tests. Their number depends only on
the fixed circumference-to-height ratio. The duality argument is
unchanged by bounded adjustments of the band edges.
\end{proof}

\subsection{Crossing inputs used before the seed}
\label{annulus:inputs}

For a rectangle $R$ contained in the cylindrical part of the ambient
sphere, write $\mathsf H_T^+(R)$ and $\mathsf V_T^+(R)$ for horizontal
and vertical crossings by occupied $B$-plus cells of class $T$.
The paths and their endpoints are vertices of the induced triangular
\emph{cell graph}.  They are not required to reach prescribed honeycomb
contour edges.  The honeycomb interpolation described after
\eqref{discrete:cell-table} will be used only after a site connection has
been obtained.

We use two crossing results from the proof of the weighted RSW estimate
in \cite[Appendix~A, pp.~46--47]{GL2025}.  Their proofs use local
comparison, site duality, and association; they precede the seed estimate
(58) in their dependency order.  In the notation $\mathrm{Cyl}_{M,N}$
of that source, the RSW estimate is
\begin{equation}\label{annulus:RSW}
 P_T\bigl(\mathsf H_T^+([-2r,2r]\times[k,k+r])\bigr)
 \geq
 \psi\!\left(P_T\bigl(\mathsf V_T^+([-3r,3r]\times[k,k+r])\bigr)\right),
\end{equation}
where $r+k\leq N$, $2r<M$, and a positive lower bound for the argument
gives a positive lower bound for the right side.  This is equation (56)
of \cite{GL2025}.  All rectangles used below are embedded, with their
horizontal length strictly less than the circumference.

The second input is horizontal lengthening: positive-probability
occupied crossings of rectangles of width four times their height can
be joined in overlapping rectangles to produce a crossing of any fixed
larger aspect ratio that fits in an embedded rectangle.  The constant
may depend on that aspect ratio.  The proof joins the uppermost crossing
in a left parallelogram to the lowermost crossing in an overlapping right
parallelogram.  We also need a conditional vertical-crossing estimate for globally
separated extremal paths.  Lemma~\ref{annulus:gluing} derives that
estimate by a stepped-domain comparison.

These assertions hold with either $T=Y$ or $T=Z$, carrying the choice
of cell class in the cylinder's boundary convention.  This is not an
identification of $P_Y$ and $P_Z$.  The RSW proof uses horizontal
translations, reflection in vertical symmetry axes, and the comparison
that moves a plus test upward.  Each of the two cell classes has those
properties.  The local reflection constructions can likewise be made
for either class.  Thus the same proof, with its chosen class fixed,
gives \eqref{annulus:RSW} and lengthening for both laws.

The weighted inputs here are the site versions specified in
\cite[Section~3.9 and Lemma~3.17]{GL2025}.  In particular they use
triangular-site duality, whose opposite crossing ends on the appropriate
\emph{site} side.  For joining paths the connector
ends on, or at an occupied neighbor of, an already occupied site path;
this is a connection in the same graph and needs no extra terminal edge.



\begin{lemma}[Finite boundary ports]\label{annulus:ports}
Consider a cell law with no prescribed complementary-spin labels.
Suppose that, after a first-spin/status history is fixed outside a
retained set, every complementary-spin connection to the exterior passes
through at most $K$ sites.  The retained cell factors, first-spin pins,
and equality constraints are kept fixed.  Removing the exterior
connections changes the marginal density of the retained $(B,\tau)$
variables by a factor between $2^{-(K-1)_+}$ and $2^{(K-1)_+}$.
\end{lemma}
\begin{proof}
For fixed retained $(B,\tau)$, the complementary-spin sum is two to the
number of components of the graph of closed cells and equality edges.
Exterior components missing all ports contribute a common factor.
The other exterior components only wire together a partition of the
ports, merging at most $(K-1)_+$ retained components.  The ratio of
unnormalized weights, after removal of the common exterior factor, is
therefore in $[2^{-(K-1)_+},1]$.  The same bounds hold for the ratio of
partition sums.  Division proves the assertion.  Constant complementary
spins satisfy all these equalities, so the two first-spin/status laws
have the same support.
\end{proof}

\begin{lemma}[Compatible finite boundary patches]
\label{annulus:patch}
Fix $x$ and a bound $M$ on the number of cells and sites of a patch.
Keep the first-spin/status history outside the patch fixed, and reveal
no complementary-spin labels.  Let $A$ be prescriptions in the patch
with the following property: for every feasible exterior configuration,
its first spins extend to the patch satisfying $A$ and all common pins,
with every prescribed open cell constant of its specified sign and every
prescribed closed cell either mixed or of positive weight $1-q_y$.
Assume that all cells incident to a mutable patch site are included in
the patch.  Then
\[
             \Prob(A\mid\hbox{exterior }(B,\tau))\geq c(x,M)>0.
\]
Consequently conditioning on $A$ changes the probability of any event
measurable off the patch by bounded factors depending only on $x,M$.
\end{lemma}
\begin{proof}
At most $M$ exterior complementary-spin components meet the patch.
Restrict their free labels to agree, losing a factor at most
$2^{(M-1)_+}$, and set all patch complementary spins to this common
label.  Choose the compatible first-spin extension.  On every
unprescribed cell select its open status if homogeneous, and its closed
status if mixed.  All its local weights are positive.  If $a(x)>0$ is
the least positive member of $x,x^2,1-x,1-x^2$, omitting zero entries,
the resulting local weight is at least $a(x)^M$.  There are at most
$2^M3^M2^M$ local spin/status assignments, each of weight at most one.
These bounds give a uniform lower bound for the ratio of the restricted
and unrestricted conditional partition sums.  In particular the
argument includes $x=1$: a prescribed closed cell must then be mixed.
For an exterior marginal $\mu$, conditioning multiplies its density by
$f/\mu(f)$, where $c\leq f=\Prob(A\mid\hbox{exterior})\leq1$.
This proves the last assertion.
\end{proof}

We will use this lemma only with a fixed-width collar between the
changed sites and any incompatible exterior pin.  Common occupied
paths prescribe their own first-spin sites and are retained throughout.
In each application below the other fixed sites are on a disjoint
minus wall.  Extend these compatible first-spin values across the
collar arbitrarily and set the complementary spin constant.  This
verifies the extension hypothesis; it does not require a finite-energy
assertion for a site already fixed by an opposite open cell.

\begin{lemma}[Equalities on a fixed plus layer]
\label{annulus:plus-equalities}
In the law of Lemma~\ref{annulus:ambient}, fix a permitted history
and an equality set $E_0$. Let $F$ be a set of edges both of whose
endpoints are pinned $B=+$. Adding the equalities $W_u=W_v$ for
$uv\in F$ stochastically decreases the joint law of $B$ and signed
$B$ statuses. There is no restriction on the number of edges in $F$.
\end{lemma}
\begin{proof}
Use $E_-=E_0\cup F$ for the lower law and $E_+=E_0$ for the upper
law in \eqref{discrete:wall-factor}. For first-spin configurations
$b,c$, put $m=b\wedge c$ and $j=b\vee c$. On $F$ the plus-wall
vectors satisfy
\[
 a_-^+(m)=0=a_-^+(b)\wedge a_+^+(c),\qquad
 a_+^+(j)\geq a_+^+(c)=a_-^+(b)\vee a_+^+(c).
\]
All minus-wall coordinates on $F$ vanish, since its endpoints are
common plus pins. Off $F$ the four inequalities in the proof of
Proposition~\ref{prop:history-order} are unchanged. The cell factors
$r_y$ do not depend on the equality set. Thus monotonicity and
log-supermodularity of $Z$ prove the same two-law inequality for the
first-spin marginals.

To restore statuses, first couple these marginals in order. In the
conditional plus-cell weights \eqref{discrete:subset-weight}, write
$A_-(S)$ and $A_+(T)$ for the allowed wall-edge vectors of lower and
upper open-cell sets. Their sorted vectors obey
\[
 A_-(S\cap T)\geq A_-(S)\wedge A_+(T),\qquad
 A_+(S\cup T)\geq A_-(S)\vee A_+(T).
\]
On $F$ the first comparison is $0\geq0$ and the second is the
inclusion of the upper open set in its union; elsewhere these are the
usual cell-set comparisons. The Bernoulli factors cancel in the
resulting two-law inequality. No cell containing an edge of $F$ can
be eligible for minus occupation, so the minus-status comparison is
unchanged. The ordered status kernels complete the joint coupling.
The nonnegative-weight argument applies also at $x=1$.
\end{proof}

\begin{lemma}[A straight short-direction crossing]
\label{annulus:straight-short}
Fix a cell class $T$ and $0<\eta<1/2$. In triangular coordinates let
$P=[0,3r]\times[0,r]$. Give its long sides open-plus $B$ cuts and
its short sides open-minus $B$ cuts, exclude their boundary-cell
factors, and leave complementary-spin labels unobserved. Use feasible
bounded transition patches at the four corners. There is
$c=c(x,\eta)>0$, independent of $r$, such that, for all sufficiently large $r$, an occupied-plus
$T$-site path in $P$ crosses from height $\eta r$ to height
$(1-\eta)r$ with probability at least $c$.
\end{lemma}
\begin{proof}
All crossings below use the induced triangular site graph. Lattice
rows and boundary contours are chosen with the fixed collars of
Lemma~\ref{annulus:patch}; that lemma absorbs the bounded corner
changes. The crossing between the two displayed interior heights is
measurable away from these patches for sufficiently large $r$.

Divide $P$ into three rhombi $R_1,R_2,R_3$ of width $r$. An ordinary
$B$-plus vertical crossing in each outer rhombus has probability bounded
below. Indeed, impose an adverse open-minus cut on its inner side.
Conditional order and the finite corner/port comparison reduce the
bound to an alternating open-arc rhombus. Reflection in its ascending
diagonal, followed by reversal of $B$, preserves that law and exchanges
an ordinary plus vertical crossing with a minus horizontal crossing.
Ordinary triangular-site duality gives a positive lower bound for
these two symmetric alternatives. Both outer crossing events are
increasing; association gives their intersection probability at least
$c_0(x)>0$. The auxiliary adverse cuts were used only to compare the
two individual probabilities and are not retained.

Run the following searches in the original law, without conditioning
on the two ordinary-crossing events.  Explore the $B$-minus component
attached to the left short side, and similarly the component attached
to the right short side.  Reveal first spins only, and declare failure
if either component reaches the middle rhombus. Retain the full Peierls interface
separating each explored component from the middle rhombus, including
all its excursions. The intersection of the two ordinary-crossing events guarantees that
each interface is confined to its outer rhombus, since its ordinary
plus crossing blocks the attached minus component.  Thus the successful
records have probability at least $c_0(x)$.  On each such record, fill
holes on the explored side and retain the region $D_0$
between the two interfaces. Its interior is unobserved, and
\[
                 R_2\subseteq D_0\subseteq P.
\]
Along either lateral interface the outer first spin is minus and the
inner first spin is plus. Each selected mixed triangle contributes the
constant $q=x^2$ and forces its three complementary spins to agree.
Successive selected triangles share a face. Thus each full lateral
interface supplies one connected $W$-equality chain, regardless of its
length or excursions. The constants cancel after conditioning on the
exploration record. Away from the bounded corner patches, the top and
bottom remain open-plus cuts and impose no $W$ equality.

Reflect in the ascending diagonal of $R_2$. In these coordinates
\[
 \rho(u,v)=(r+v,u-r),\qquad D=D_0\cup\rho(D_0).
\]
This reflection preserves the chosen cell class. For example, with
$\omega=e^{i\pi/3}$, $\Lambda=\Z+\omega\Z$ and
$Y=\Lambda-i/\sqrt3$, its linear part is $z\mapsto\omega\bar z$,
and $\omega\overline{(-i/\sqrt3)}+i/\sqrt3=\omega-1\in\Lambda$.
The same holds for $Z$. The inclusions above give
\[
 \rho(D_0)\cap P=R_2,\qquad D\cap P=D_0.
\]
The union is a simply connected symmetric domain: its lower and upper
reflected arms attach to $D_0$ along the bottom and top sides of $R_2$.
The construction is made on the actual lattice contours with matching
endpoint collars.

Put $A=(r,0)$, $C=(2r,r)$, $B_0=(2r,0)$ and $D_1=(r,r)$.
Reflection fixes $A,C$ and exchanges $B_0,D_1$. The boundary arc
$AB_0$ surrounds the lower reflected arm, and $CD_1$ surrounds the
upper one. Any occupied site connection between these two arcs has a
bottom-to-top subcrossing of $P$ entirely in $D_0$: take its segment
after the last bottom-side visit before the first subsequent top-side
visit. It therefore contains the crossing between the two interior
heights in the statement. No confinement of this subcrossing to $R_2$
is required.

On $D$ take the excluded-boundary-cell law with $B=+$ on its inner
boundary layer and with $W$ constant, of unspecified sign, on that
layer. Denote it by $\nu_D$. This is the equal mixture of the two laws
with boundary spin pairs $(+,+)$ and $(+,-)$; global $W$ reversal
preserves all $B$-occupation events. For the $(+,+)$ law, occupied-site
super-duality, reflection and interchange of colors give
\begin{equation}\label{annulus:symmetric-short}
 \nu_D(AB_0\mathrel{\longleftrightarrow}_{\xi^{B,+}}CD_1)
 \geq\tfrac12.
\end{equation}
This is precisely the site-arc convention of
\cite[Lemma~3.17]{GL2025}: the two alternative color crossings have
union of probability one, and symmetry equates their probabilities.
Their disjointness is not required. The sign of a crossing incident to
the corresponding boundary layer is plus.

Where the old top and bottom cuts lie in the interior of $D$, prescribe
them open-plus. These are increasing first-spin/status prescriptions,
so association preserves the lower bound
\eqref{annulus:symmetric-short}. Their open factors are constants and
remove all $W$ interaction across them. On old top/bottom pieces that
lie on $\partial D$, the boundary factors are already excluded; no
additional exterior minus pins are imposed on the sites of those
open cells. This convention avoids a conflict between an open-plus
cell and an exterior minus spin in the same cell.

We compare the resulting law in $D_0$ with the law left by the actual
exploration. The old top/bottom pieces outside the middle rhombus lie
on $\partial D$ and hence have constant $W$ under $\nu_D$. In the
actual law add exactly these equalities, along edges whose endpoints
are already pinned $B=+$ by the open cut. By
Lemma~\ref{annulus:plus-equalities}, this decreases all increasing
$B$-occupation events, with no loss depending on the number of edges.
No equality is added on the old top/bottom pieces that run through the
interior of $D$, namely the middle-rhombus pieces.

After this lowering, the two laws have identical retained cell
factors, first-spin pins, and long boundary equality chains. The two
lateral chains, extended along their adjacent old top/bottom boundary
pieces, contain every remaining complementary-spin connection to the
exterior, except for the bounded transition patches. Their sites all
have fixed boundary first-spin data. The remaining exterior can only
identify these two chains or leave their labels separate. For fixed
retained $(B,\tau)$, such an identification changes the complementary
spin sum by a factor in $\{1,1/2\}$; after normalization the two
marginal densities differ by at most a factor two. Mixtures over
unobserved exterior statuses preserve this bound. Applying
Lemmas~\ref{annulus:ports} and \ref{annulus:patch} to the finitely many
additional corner components gives a uniform factor $c_1(x)>0$.

Consequently, for every successful ordinary-interface exploration
record, the original probability of the occupied crossing between the
interior heights is at least $c_1(x)/2$. Averaging over those records
and using their probability at least $c_0(x)$ proves the assertion.
All modifications at the corners are compatible bounded patches and
all other comparisons are monotone or involve the two connected
boundary chains. Thus no bound depends on a boundary length. The
argument and the patch lemma include $x=1$.
\end{proof}

\begin{figure}[tbp]
 \centering
 \begin{tikzpicture}[x=0.9cm,y=3.0cm,>=Latex,font=\small]
  \definecolor{annulusplus}{RGB}{25,90,150}
  \definecolor{annulusminus}{RGB}{150,65,45}
  \fill[black!4] (0,1.25)--(6,1.25)--(6,.5)--(7,.5)--(7,-.25)--(1,-.25)--(1,.5)--(0,.5)--cycle;
  \fill[annulusplus!10] (1.25,-.25) rectangle (5.75,1.25);
  \fill[annulusplus!22] (1.25,.3) rectangle (5.75,.7);
  \draw[annulusplus,very thick] (0,1.25)--(6,1.25) node[midway,above] {$+$};
  \draw[annulusplus,very thick] (1,-.25)--(7,-.25) node[midway,below] {$+$};
  \draw[annulusminus,thick] (0,1.25)--(0,.5)--(1,.5)--(1,-.25);
  \draw[annulusminus,thick] (6,1.25)--(6,.5)--(7,.5)--(7,-.25);
  \node[annulusminus,left] at (0,.58) {$-$};
  \node[annulusminus,right] at (7,.42) {$-$};
  \draw[annulusplus,thick] (0,.94) .. controls (.8,.86) and (1.3,.98) .. (2,.89)
    .. controls (2.8,.79) and (3.1,.96) .. (4,.87)
    .. controls (4.9,.79) and (5.2,.98) .. (6,.9);
  \draw[annulusplus,thick] (1,.1) .. controls (1.8,.17) and (2.3,.05) .. (3,.13)
    .. controls (3.8,.22) and (4.1,.04) .. (5,.15)
    .. controls (5.8,.21) and (6.3,.02) .. (7,.11);
  \node[annulusplus,above] at (3.4,.99) {$\Gamma_L$};
  \node[annulusplus,below] at (4,.05) {$\Gamma_R$};
  \foreach \u/\v in {0/.94,6/.9,1/.1,7/.11}{
    \fill[orange!65] (\u,\v) ellipse [x radius=.11,y radius=.035];
  }
  \draw[dashed,black!60] (1.25,-.25)--(1.25,1.25);
  \draw[dashed,black!60] (5.75,-.25)--(5.75,1.25);
  \node at (3.5,.5) {$E:\ \text{central vertical crossing}$};
  \node[below] at (1.25,-.28) {$5\ell/4$};
  \node[below] at (5.75,-.28) {$23\ell/4$};
  \node[right] at (6.1,1.25) {$5\ell/4$};
  \node[right] at (7.1,-.25) {$-\ell/4$};
  \node[above left] at (0,1.25) {$S$};
\end{tikzpicture}

 \caption{The stepped reference region for the separated-path comparison.
 The full upper and lower paths stay outside the central crossing band.
 Orange patches contain the four actual endpoint records.  The dashed
 sides bound the central $3:1$ comparison parallelogram; both those
 sides and the event $E$ are separated from the endpoint patches.
 The vertical coordinate is expanded, and fixed-width outward shifts
 and corner rounding are omitted.}
 \label{annulus:fig-stepped}
\end{figure}

\begin{lemma}[A separated-path crossing]
\label{annulus:gluing}
Use triangular coordinates $(u,v)$, with $u$ horizontal, and a fixed
cell class $T$.  Let
\[
 P_L=[0,6\ell]\times[0,\ell],\qquad
 P_R=[\ell,7\ell]\times[0,\ell]
\]
be embedded parallelograms in a matching mixed cylinder.  Let $\Gamma_L$
be the uppermost occupied-plus horizontal crossing of $P_L$ and
$\Gamma_R$ the lowermost one of $P_R$.  Suppose that $\Gamma_L$ lies
above $3\ell/4$ and $\Gamma_R$ below $\ell/4$.
For all sufficiently large $\ell$ there is $c=c(x)>0$, uniform in
the feasible extremal pair and in $\ell$, such that, conditional on
this pair, with probability at least
$c$ there is an occupied-plus vertical crossing between heights
$3\ell/10$ and $7\ell/10$ in the horizontal window
$[5\ell/4,23\ell/4]$.
The cylinder and lattice rounding are required to leave a fixed-width
collar around the displayed parallelograms.
\end{lemma}
\begin{proof}
We prove the estimate uniformly after refining the two extremal paths
to their exterior exploration records.  The upper exploration reveals
only the part of $P_L$ above its full crossing, and the lower one only
the part of $P_R$ below its full crossing, together with the prescribed
first-spin layers on the crossings.  All observed statuses include
all their first-spin sites.  No complementary-spin label is revealed.
The path version of the invasion argument in
Lemma~\ref{annulus:search} gives these records.  Use the actual
authorized honeycomb boundary of the explored face set to realize each
extremal crossing.  Every queried cell meeting both sides is then an
open boundary cell; its retained first spins are its prescribed plus
layer.  This avoids revealing an additional strip along an arbitrary
interpolation of a site path.  We retain the full crossings, including
all their excursions near a side.

Construct a stepped reference region $S$.  Its upper portion has
horizontal range $[0,6\ell]$ and vertical range
$[\ell/2,5\ell/4]$; its lower portion has horizontal range
$[\ell,7\ell]$ and vertical range $[-\ell/4,\ell/2]$.
Move each of its two stepped sidewalls outward by a fixed number of
lattice spacings, say twenty, and round to honeycomb paths.  The top
and bottom sides are likewise chosen as actual lattice contours.
The steps lie at height near $\ell/2$, at positive macroscopic distance
from both extremal paths.

Join the actual four endpoints of the full crossings to the respective
outer sidewalls through four bounded patches.  The region $D$ between
the two full paths and these sidewalls is unobserved outside those
patches.  Indeed its upper part is below the full $\Gamma_L$ in $P_L$,
and its lower part is above the full $\Gamma_R$ in $P_R$.  In their
overlap the two statements hold simultaneously, since the height bands
are separated.  The narrow strips introduced outside a side of a
parallelogram were not queried by its exploration; the other exploration
is separated from them in height.  Thus the only extra recorded sites
or cells introduced by an endpoint join lie in its bounded patch.
This uses the full paths: no discarded prefix is placed in $D$.

Prescribe the two adverse sidewalls open-minus, stopping them before
the endpoint patches.  They lie in the unobserved region and are
disjoint from the occupied-plus crossings.  By conditional order this
can only decrease the probability of the desired central vertical
crossing; the rarity of the prescribed walls incurs no probability
factor.  We next make bounded, compatible completions in the endpoint
patches.  Let $h$ be the complete original record in those patches:
retain every recorded first spin and status, including all first-spin
sites of each recorded cell.  Retain also the plus pins forced by the
crossings.  The new minus walls are separated from these sites by the
fixed outward displacement.  Choose a compatible completion $A$ only
in a bounded core, keeping an outer collar unprescribed except for the
common crossing and wall pins.  The fixed first-spin values are
compatible; extend them across this free collar and set the
complementary spin constant.  On unprescribed core cells choose open
status when homogeneous and closed status when mixed.  Include every
cell incident to a mutable core site in the patch.  Thus
Lemma~\ref{annulus:patch} gives a uniform positive lower bound for $A$.
At $x=1$, recorded closed cells keep their already recorded mixed first
spins.  Write $h^*=(h,A)$, and keep this complete bounded history in all
reference laws below.  No arbitrary part of the original record is
removed by a finite-energy assertion.

Let $\nu_{S,h^*}$ be a reference law with open-plus top and bottom arcs,
open-minus stepped sidewalls, and this same history $h^*$ in the four
endpoint cores.  Realize the law in a filled ambient triangulation by
these cell and first-spin prescriptions, with a fixed feasible
completion of its outer corner cells and no complementary-spin labels.
Its top and bottom are at distance at least $\ell/4$ from the original
parallelograms, so their prescribed plus layers do not meet $h^*$.
The law is positively associated by Lemma~\ref{annulus:ambient}.

Condition this reference law on the two full paths being occupied-plus.
The law in $D$ has the same retained cell factors and first-spin pins
as the original record completed by $A$.  Every other boundary cell is
open and removes the complementary-spin interaction across its arc.
Choose each completed core to contain its entire short endpoint join,
including every non-open cell crossing the chosen boundary of $D$ and
all its first-spin sites.  The boundary enters and leaves the free outer
collar only along the common open crossing or wall arcs.  Thus every
cross-boundary cell outside the completed cores is already open.  All
other cross-boundary first spins and statuses are fixed in $h^*$; only
the uniformly bounded number of complementary-spin sites in those
cores can communicate with the exterior.  Lemma~\ref{annulus:ports} compares the
two retained laws with a uniform factor, independently of the exterior
complementary-spin wiring.  The original endpoint history $h^*$ is
kept on both sides of this comparison.

Write $E$ for the vertical crossing of the central height interval in
the statement.  Its support, including the incident first-spin sites of its cells,
lies in $D$ for sufficiently large $\ell$; the full comparison rectangle
need not lie in $D$.  The bounded displacement between the occupied site
path and its authorized exploration boundary is absorbed by the strict
height gaps.  Its horizontal window is at macroscopic distance from
all actual endpoint patches, whose horizontal coordinates are near
$0,6\ell,\ell,7\ell$.  The fixed corners of $S$ are also outside this
window.  The lower bound for $A$, followed by the port comparison,
gives
\begin{equation}\label{annulus:record-comparison}
 \Prob(E\mid\hbox{original record})
 \geq c_1(x)\,
 \nu_{S,h^*}(E\mid\Gamma_L,\Gamma_R\hbox{ occupied-plus})
 \geq c_1(x)\,\nu_{S,h^*}(E).
\end{equation}
The last inequality is association: the full-path prescriptions are
increasing, with every first-spin site of each path cell included.
All constants are uniform in the retained bounded history $h^*$.

It remains to bound $\nu_{S,h^*}(E)$.  In the reference law alone, prescribe
two additional open-minus walls at horizontal positions $5\ell/4$
and $23\ell/4$, ending a bounded distance before the top and bottom
plus arcs.  This decreases $E$.  Complete the four new transition
patches as above.  These walls and patches are at macroscopic distance
from every actual endpoint patch.  The central region is a straight
parallelogram of width $9\ell/2$ and height $3\ell/2$, hence aspect
ratio three.  Exact cut factorization off its corners and
Lemmas~\ref{annulus:ports}--\ref{annulus:patch} compare its law with
the alternating open-arc law of
Lemma~\ref{annulus:straight-short}.  Lemma~\ref{annulus:straight-short}, with $r=3\ell/2$ and
$\eta=11/30$, bounds below the probability of the crossing $E$ itself:
its two heights are
$-\ell/4+(11/30)(3\ell/2)=3\ell/10$ and
$-\ell/4+(19/30)(3\ell/2)=7\ell/10$.
It follows that
$\nu_{S,h^*}(E)\geq c_2(x)>0$.

The additional interior minus walls were imposed only in $\nu_{S,h^*}$,
after removing the conditioning on the actual crossings in
\eqref{annulus:record-comparison}.  They are therefore not required to
be compatible with paths that might cross them.  Combining the two
bounds and averaging over the exterior exploration records proves the
lemma.  All exceptional corner adjustments involve a bounded number
of cells, with free outer collars and compatible common pins as above.
\end{proof}

\subsection{A seed with either cell orientation}

We give the scale and extraction conventions once.  Fix $C\geq3$ and
work in the mixed cylinder in (58) of \cite{GL2025}, with circumference
$L\sim2Cn$ and height $H\sim5n$.  These relations only allow the usual
bounded changes needed to select actual lattice rows.  Choose
$m=\lfloor\varepsilon H\rfloor$, where $\varepsilon>0$ is a fixed small
constant, for instance $10^{-5}$.  Every construction below fits with
this choice; making $\varepsilon$ smaller by a fixed factor would only
change constants.  We first take $n$ large enough for all clearances to
exceed a fixed number of lattice spacings.

A band of height $m+O(1)$ supplies a finite family of tests as follows.
Lift paths to the flat cover of the cylinder, and put horizontal centers
on a net of spacing at most $m/4$.  If a vertical crossing has horizontal
extent at most $5m$, a window of width $5.5m$ contains it.  Trim its
vertical ends to obtain a vertical test of width six times its height,
with height between $.96m$ and $.99m$, after the allowed lattice
rounding.  If its extent is greater than $5m$, take a subpath between
vertical lines separated by $4.4m$ and trim it to a horizontal test
of width four times its height, with height between $1.01m$ and
$1.05m$.  The latter window encloses the original band vertically.
An occupied winding also supplies horizontal tests of this type.  The
slack permits the specified centers, lattice rows, and crossing-side
conventions to be adjusted before trimming.  There are finitely many
tests, with their number bounded in terms of $C$ and $\varepsilon$, not
$n$.  We call a lower bound for one of these tests a \emph{seed}.

\begin{lemma}[Seed with a matching class]\label{annulus:first-seed}
There is $c=c(x,C)>0$ such that, for each sufficiently large $n$, one of
$T=Y,Z$ has a seed of occupied $B$-plus cells, under the matching mixed
law $P_T$, in a band of height comparable to $m$ centered near $.71H$.
The seed is a horizontal crossing of aspect ratio four or a vertical
crossing of aspect ratio six.  The choice of class and test may depend
on $n$; the lower bound $c$ does not.
\end{lemma}
\begin{proof}
Start under $P_Y$ with a thin band $S$ centered near $H/3$.
Lemma~\ref{annulus:duality} gives a vertical occupied $B$ crossing of
one sign, or a winding occupied $W$ circuit of one sign.  The sum of
the probabilities of these four events is at least one.

A $B$-plus vertical crossing gives a $Y$ seed by the finite-family
extraction above.  The slide comparison of Lemma~\ref{annulus:slide}
moves its test upward.  A $B$-minus vertical crossing is reflected
across the horizontal middle line of the cylinder, followed by a flip
of $B$.  The resulting law is $P_Z$, including its boundary-cell
convention: horizontal reflection interchanges the two cell classes.
It gives a $Z$ plus seed near $2H/3$, which can again be moved upward.
Neither step asserts a color-preserving symmetry of a fixed $P_Y$ law.

It remains to treat a $W$ winding.  Global $W$ flip lets us take its
sign to be plus, at a loss of at most a factor two.  Search from below
in a slightly enlarged band, as in Lemma~\ref{annulus:search}.  On a
successful record the retained upper region lies above one open
$W$-plus winding $\gamma$ of class $Y$ and below the original upper
$B$-plus cut.  All tests used next lie in its unobserved part.

Choose a horizontal lattice reflection axis near $.69H$.  The reflection
of $\gamma$ lies strictly above the old top, since
$2(.69)-1/3>1$ and the search band is thin.  Move the top $B$-plus cut
up to that reflected contour, now of class $Z$.  By
Lemma~\ref{annulus:one-white} and the cut comparison, this move
\emph{decreases} every increasing $B$-plus test.  The law between
$\gamma$ and its reflection is invariant under reflection followed by
exchange of the two spin colors.  This follows from the coherent edge
weight and the exact cut factors; the transformation also transports
$Y$ tests to $Z$ tests.

Apply band duality near the reflection axis, using fresh $Y$
augmentation.  The finite extraction gives a family of $B$ seeds of
either sign and $W$ horizontal seeds of either sign whose probabilities
have sum bounded below.  Reflection and color exchange turn the $W$
seeds into $Z$-oriented $B$ seeds.  For either orientation a minus seed
has probability at most that of the corresponding plus seed: flipping
$B$ changes the top prescribed cut from plus to minus while retaining
the one $W$-plus cut, and Lemma~\ref{annulus:one-white} orders these
two laws.  Consequently a fixed finite family of $B$-plus tests, in the
two orientations, has total probability bounded below in the symmetric
law, uniformly over the successful search record.

The same lower bound holds in the original law above $\gamma$, by
the comparison that moved the top upward.  Its index family can be
chosen independently of $\gamma$: the search band, reflection axis,
and enlarged test windows are deterministic and separated.  Averaging
over successful records therefore gives a positive lower bound for
one of those tests under $P_Y$.  If that test uses class $Z$, its
augmentation is fresh and unobserved.  Apply the slide comparison,
changing the boundary class to that of the test, to move it upward and
obtain an ordinary matching $P_Z$ experiment.  The same comparison
moves a $Y$ test upward within $P_Y$.

All three branches can thus be placed near $.71H$, with clearance
between every old and new cut when their classes differ.  Taking the
minimum of the finitely many bounds proves the lemma.  The sum bound is
obtained before choosing its successful member, so no random test is
selected inside a conditional probability.
\end{proof}

\subsection{A winding and recovery of the requested orientation}

Apply \eqref{annulus:RSW} if the seed in
Lemma~\ref{annulus:first-seed} is vertical.  We have a positive lower
bound for an occupied horizontal plus crossing in the matching law
$P_T$.  If $T=Y$, horizontal lengthening gives the horizontal crossing
required in (58), inside its target band $[3n,4n]$: the seed is near
$.71H\sim3.55n$ and its height is much smaller than $n$.  It remains
to treat $T=Z$.  We do not identify opposite-orientation augmentations
at the same location.

\begin{lemma}[Winding from the matching seed]\label{annulus:winding}
Under $P_Z$, the horizontal seed above implies a positive lower bound
for an occupied $B$-plus winding of class $Z$ in a thin neighborhood of
height $.71H$, entirely below $.73H$.  The bound is uniform in $n$.
\end{lemma}
\begin{proof}
We first obtain a short vertical crossing in a controlled horizontal
window.  Use the two overlapping width-$6\ell$ parallelograms of
Lemma~\ref{annulus:gluing}, choosing their vertical height to be about
$50m$.  Their total horizontal span is less than $500m$.  Place their bottom at, or just
above, the seed band.  In a thin band above relative height $40m$,
request a horizontal crossing through the left parallelogram; in one
below relative height $10m$, request a horizontal crossing through the
right parallelogram.  These estimates follow from the original
horizontal seed, upward translation by Lemma~\ref{annulus:slide}, and
horizontal lengthening with fixed aspect ratios.  Cross rectangles
extending beyond the parallelogram's sloping sides and then trim the
paths to those sides.  Association makes the intersection of the two
requests have probability bounded below.

Let $\Gamma_L$ be the uppermost crossing of the full left
parallelogram and $\Gamma_R$ the lowermost crossing of the full right
one.  A crossing in the upper thin band separates the top side from the
region below that band; hence the uppermost crossing stays above
relative height $39m$, allowing lattice rounding.  Similarly the
lowermost right crossing stays below height $11m$.  Thus the two
thin-band requests imply that the pair $(\Gamma_L,\Gamma_R)$ lies in
a set $\mathcal G$ of separated realizations in gluing order.

For each realization in $\mathcal G$,
Lemma~\ref{annulus:gluing} gives a vertical crossing of the central gap
with conditional probability at least $c_0>0$.  The bounds $39m$ and
$11m$ put the extremals above three quarters and below one quarter of
their height-$50m$ parallelograms.  We condition only on the two
extremal crossings, not additionally on the thin-band requests.
Averaging gives a lower bound $c_0\Prob(\mathcal G)$, and
$\Prob(\mathcal G)$ is at least the probability of the requests.
The resulting crossing runs between relative heights $15m$ and $35m$
in a horizontal window of width at most $500m$.  This supplies the
required local vertical-crossing estimate.

Place a thin horizontal band about relative height $25m$, strictly
between these two vertical levels.  Translate the vertical-crossing
window to two seam neighborhoods separated by about $L/2$.  Their
widths are much smaller than $L$.  Request also two horizontal
crossings in the thin band: one follows each of the two arcs between
the seams, and each extends beyond the whole vertical window at both
ends.  These horizontal rectangles have length, for example, about
$3L/4$, so each is embedded in the cylinder.  Their crossing estimates
are ordinary horizontal lengthenings.  All four events are increasing
in the same signed $B$-status field.  Association bounds the probability
of their intersection below by a positive constant.

On that intersection both horizontal paths meet the vertical path in
each seam neighborhood, by planar site crossing.  In the flat cover,
the first horizontal path connects seam 1 to seam 2 on its right, and
the second connects seam 2 to the next periodic copy of seam 1.  The
connections therefore join a site path to its translate by one full
period.  Its projection contains a noncontractible occupied circuit.
The interpolation to authorized honeycomb edges changes its position
by only a bounded number of lattice spacings and yields a simple
winding cut in a slightly larger band.  All these bands remain below
$.73H$ because $m/H$ is fixed and sufficiently small.
\end{proof}

\begin{proposition}[The mixed-cylinder seed estimate]
\label{annulus:seed}
For every fixed $C\geq3$ there is $\delta=\delta(x,C)>0$ such that,
for all sufficiently large $n$,
\begin{equation}\label{annulus:seed-bound}
 P_Y\bigl(\mathsf H_Y^+([-2n,2n]\times[3n,4n])\bigr)\geq\delta
\end{equation}
on the mixed cylinder $\mathrm{Cyl}_{Cn,5n}$ of
\cite[Appendix~A]{GL2025}, with the same lattice rounding convention.
\end{proposition}
\begin{proof}
Only the $T=Z$ branch remains after Lemma~\ref{annulus:first-seed}
and \eqref{annulus:RSW}.  By Lemma~\ref{annulus:winding}, with
probability bounded below an exterior search finds a $Z$-oriented
$B$-plus winding below $.73H$.  Above a successful record the retained
law is between that plus cut and the original upper plus cut.  It is
therefore above the free capped-sphere law for increasing $B$-plus
tests: prescribing either plus cut is an increase by
Lemma~\ref{annulus:ambient}.  Tests of class $Y$ well above the
explored band use disjoint, fresh cells, so this comparison applies to
them too.

Test a thin band near $.75H$ using class $Y$.  The free sphere law is
invariant under the two spin flips and exchange of the spin colors.
Band duality and finite extraction therefore imply a positive lower
bound for the sum of a fixed finite family of $Y$-oriented $B$-plus
seed probabilities.  These spin symmetries do not require spatial
translation or reflection symmetry of the caps.  The lower bound holds
in every retained law between the two plus cuts.  Average over
successful records, and then choose one member of the deterministic
family.  This gives a $Y$ seed under $P_Z$.

Finally slide the boundary cuts and change their class to $Y$, moving
the test to height near $.77H$.  Lemma~\ref{annulus:slide} preserves
the lower bound and gives the matching mixed law $P_Y$.  The clearance
between old and new cuts guarantees that the cell families used in
this comparison are edge-disjoint.  Apply \eqref{annulus:RSW} if the
seed is vertical, then horizontally lengthen to width $4n$.  Since
$.77H\sim3.85n$ and all vertical bands used here are much thinner than
$n$, the resulting path stays in $[3n,4n]$.  It is the crossing in
\eqref{annulus:seed-bound}.
\end{proof}

\subsection{The annulus estimate and the field theorem}

\begin{proof}[Proof of Theorem~\ref{thm:circuit-input}]
Proposition~\ref{annulus:seed} proves the horizontal alternative in
(58) of \cite{GL2025}.  We now use the implications of that source's
Appendix~A.  The RSW estimate (56) and horizontal lengthening convert
(58) into the pushing estimate (57); the comparison between rectangle
and cylinder has exactly the boundary-cell convention realized by our
cuts.  Pushing gives the renormalization inequality (23), and its
non-exponential-decay branch, selected by the exclusion argument on
pp.~44--45 of \cite{GL2025}, gives equation (55).  Those implications
use association, domain Markov comparison, the weighted site-crossing
inputs already specified, and qualitative delocalization.  Their output is precisely
an occupied plus circuit under an excluded-boundary-cell law with the
first color minus on its boundary layer and the second color free.
Lemma~\ref{lem:cell-cut} identifies this with the law used in the
statement of Theorem~\ref{thm:circuit-input}.

Our seed proof was for sufficiently large lattice scales.  The finitely
many remaining admissible $n$ can be absorbed by decreasing the
constant.  The condition $n\geq3/(a-1)$ implies
$\lfloor an\rfloor-n\geq3$, so a plus ring can be placed around layer
$n+1$, with plus first spins through layer $n+2$, minus first spins on
the outer prescribed boundary layer, and constant second spin.  Force
the finitely many ring cells open.  This is a feasible positive-weight
configuration for every $x\in[1/\sqrt2,1]$.  Thus each of those finite
probabilities is positive.  Spin flips and color exchange give the
other versions of the theorem.
\end{proof}

The dependence of the proof is
\[
 \text{cut and order lemmas, (56), local site gluing}
 \ \Longrightarrow\ \eqref{annulus:seed-bound}
 \ \Longrightarrow\ (57),(23),(55).
\]
The seed uses only the local cut, order, and crossing inputs displayed
above.  Theorem~\ref{thm:circuit-input} supplies the annulus estimate
in Lemma~\ref{lem:buffer-circuits}.
Proposition~\ref{prop:plane-covariance} gives $v\in(0,\infty)$ with
the Euclidean-area convention.  Section~\ref{sec:limit} then proves
tightness, the Wick recursion, and all mixed-moment convergence with
$\sigma(x)=\sqrt v$ for the inside approximations of
Theorem~\ref{thm:main}.
